\documentclass[10pt,a4paper]{amsart}
\usepackage[margin=19mm]{geometry}
\usepackage{amsmath,amssymb,mathtools}
\usepackage[scr=rsfs,cal=euler]{mathalfa}
\usepackage{xfrac}
\usepackage[unicode,hidelinks,bookmarks=false]{hyperref}
\newcommand{\ie}{i.e.\ }
\newcommand{\cf}{cf.\ }
\newcommand{\resp}{respectively\ }
\newcommand{\Subsection}{\subsection}
\providecommand{\nobreakdash}{\nobreak\hbox{-}}
\setlength{\emergencystretch}{2em}
\multlinegap0pt
\usepackage{xcolor}
\usepackage{mathtools}
\newtheorem{lemma}{Lemma}
\newcommand{\lnn}{\ell^2(\mathbb{N})}
\title{Infinite dimensional generative sensing}
\author{Paolo Angella, Vito Paolo Pastore, Matteo Santacesaria}
\date{Source: arXiv:2603.03196v1 (CC BY 4.0)}
\begin{document}
\maketitle

\noindent\emph{Source and licence.} Infinite dimensional generative sensing. Version arXiv:2603.03196v1 (CC BY 4.0), distributed by arXiv under the Creative Commons Attribution 4.0 International licence, http://creativecommons.org/licenses/by/4.0/, as recorded in the arXiv licence metadata for this exact version and shown on the abstract page https://arxiv.org/abs/2603.03196v1. Full licence notice: \texttt{licenses/arxiv-2603.03196-CC-BY-4.0.txt}.\par\smallskip

\noindent\textit{Editorial scope.} Counted unit: the article's lemma coherence-energy-bound (source0.tex lines 519--525). The article's complete original proof of it is delivered with it (source0.tex lines 527--546, one \texttt{proof} environment of 20 lines). No mathematics is written, repaired or re-derived: the statement and the proof are the article's own lines, in the article's own notation, and no retained line is substituted or re-flowed.  No further source-local context is needed: every cross-reference of every retained line resolves inside the two delivered spans.  Nothing was omitted from any retained span, so the package carries no omission record.  No retained line contains a citation, so the wrapper carries no bibliography and no bibliography entry.  No retained line contains a figure environment, an image-inclusion command or drawing code, so no image is delivered and the package declares no asset dependency.  Editorial formatting only, in addition: an AMS \texttt{amsart} preamble that restates the article's own declarations and macros used by the retained lines, namely the environments \texttt{lemma} on separate counters, together with the standard packages those lines need.  The retained environments print in this document as Lemma 1 in order of appearance, whereas the article prints the same items with the numbering of its own full document.  The complete native source of the article as distributed in the arXiv e-print for this exact version is bundled unchanged as \texttt{sources/195661/source0.tex}.
\begin{lemma}\label{coherence-energy-bound}
    Given $F$ unitary in $\lnn$, $\mathcal{U}\subseteq\lnn$ a finite union of convex cones contained in a finite-dimensional subspace, and a probability distribution $p$ with $\hat{I}$ the projection on the support of $p$. If we have 
    \begin{equation}\label{magg-1}
        \| \hat{I} \alpha \| \geq 1,
    \end{equation}
    then $\mu_{\mathcal{U}}(F,p) \geq 1$.
\end{lemma}

\begin{proof}
  
From the definition of $ \mu_{\mathcal{U}}(F,p) $, if we denote with $\Omega$ the support of $p$ we have
    \begin{equation*}
        \alpha_j^2 \leq \mu_{\mathcal{U}}(F,p)^2 p_j \quad \forall j \in \Omega.
    \end{equation*}
    Summing this inequality over all indices in $\Omega$:
    \begin{equation*}
        \sum_{j \in \Omega} \alpha_j^2 \leq \sum_{j \in \Omega} \mu_{\mathcal{U}}(F,p)^2 p_j = \mu_{\mathcal{U}}(F,p)^2 \sum_{j \in \Omega} p_j.
    \end{equation*}
    Therefore,
    \begin{equation*}
        \sum_{j \in \Omega} \alpha_j^2 \leq \mu_{\mathcal{U}}(F,p)^2.
    \end{equation*}
    Taking the square root implies
    \begin{equation*}
        \mu_{\mathcal{U}}(F,p) \geq \sqrt{\sum_{j \in \Omega} \alpha_j^2}.
    \end{equation*}
    Thus, if \eqref{magg-1} holds, it follows directly that $\mu_{\mathcal{U}}(F,p) \geq 1$.
\end{proof}

\end{document}
